\subsection*{Local Search}
\label{sec:local-search}

The local search (LS) acquisition function optimizer in the \texttt{bbotk} package works similarly to the LS implementation in SMAC\cite{lindauer2022} and can be used to optimize acquisition functions in mixed spaces.
%Therefore, it is described here in more detail.
To reduce function call overhead, $k$ LS runs are generated concurrently, so that as many points as possible can be passed to the SM in a single batch.

Each LS run is executed for a specified number of iterations.
In each iteration, the LS generates multiple candidate neighbors by mutating the current point / incumbent.
A neighbor is obtained by mutating exactly one parameter from the current point.
For numeric parameters, the parameter value is first scaled to $[0, 1]$, zero-mean Gaussian noise (with standard deviation governed by a control parameter) is added, and the result is then rescaled to the original domain.
Afterwards, the value is clipped to the corresponding lower and upper bounds.
For integer parameters, the same procedure is applied, but the result is rounded to the nearest integer.
For factor parameters, a new level is sampled uniformly from the levels that differ from the current one.
For logical parameters, the value is simply flipped.

Hierarchical dependencies in the search space are handled as follows.
Only active parameters are eligible for mutation.
After a mutation occurs, all conditions in the search space are checked in topological order.
If a condition is not satisfied, the corresponding parameter is marked as inactive.
Conversely, if all conditions for a parameter are satisfied but it is marked as inactive, it is assigned a random valid value.

After the neighbors are generated, they are evaluated using the acquisition function.
The number of candidates passed to the SM simultaneously is therefore $k$ times the number of neighbors generated per LS run. 
For each LS, the best neighbor is selected as the new incumbent unless all neighbors are worse than the current incumbent, in which case the incumbent is retained.
To avoid getting stuck in local optima, the algorithm employs a restart mechanism: 
For each LS, the number of consecutive iterations without improvement is tracked, and if this 
exceeds a specified threshold, the search is restarted from a new random point in the search space.